\section{Introduction}

The predictability of stock returns need not be constant across market environments. Periods of elevated market volatility can coincide with changes in risk premia, investor behavior, information processing, liquidity, and the cross-sectional dispersion of returns. This research therefore asks whether next-month stock-return predictability differs between high- and low-volatility regimes and whether a nonlinear machine-learning model captures regime-dependent predictive structure differently from a regularized linear model.

The production empirical design compares Elastic Net and XGBoost using the same three-predictor information set and the same strictly out-of-sample evaluation structure. At formation month $t$, log market equity, book-to-market, momentum 12--2, and the volatility regime are constructed using information available through $t$, while the prediction target is the realized stock return in calendar month $t+1$. Models are estimated only on historical observations whose outcomes are observable by the formation date.

The production panel combines CRSP monthly stock returns with Compustat annual accounting information through the CRSP/Compustat Merged database. Annual book-to-market is assigned using a conservative June availability convention and prior-December market equity. Volatility regimes are based on VIX relative to its expanding historical median, preventing future volatility information from entering the state classification.

Hyperparameters are selected once using only the 1990--1999 pre-OOS period with forward-chaining validation and are frozen thereafter. The OOS evaluation begins in January 2000 and extends through November 2025. This design separates historical model selection from subsequent forecast evaluation and avoids random train/test splits that would violate the temporal structure of the problem.

The production results reveal an important distinction between unconditional forecast accuracy and cross-sectional economic value. Both models have OOS $R^2$ values close to zero relative to a pooled expanding historical-mean benchmark. However, XGBoost displays statistically positive cross-sectional rank information in LOW-VIX months, and equal-weighted prediction-sorted long--short portfolios retain positive returns after the frozen transaction-cost adjustment. The same signal is weak in HIGH-VIX months and disappears under value weighting. At the same time, direct LOW-minus-HIGH tests do not reach conventional significance. The evidence therefore points to predictive content that is concentrated in low-volatility states without establishing a sharp statistically significant break between regimes.
